\documentclass[12pt,a4paper]{article}

% ── Packages ────────────────────────────────────────────────────────────────
\usepackage[T1]{fontenc}
\usepackage[utf8]{inputenc}
\usepackage{lmodern}
\usepackage{microtype}
\usepackage{amsmath,amssymb,amsthm}
\usepackage{mathtools}
\usepackage{enumitem}
\usepackage{geometry}
\usepackage{hyperref}
\usepackage{booktabs}
\usepackage{array}
\usepackage{xcolor}
\usepackage{listings}
\usepackage{fancyhdr}

\geometry{margin=2.5cm}
\setlength{\headheight}{13.60001pt}
\addtolength{\topmargin}{-1.60001pt}

% ── Theorem environments ─────────────────────────────────────────────────────
\newtheorem{theorem}{Theorem}[section]
\newtheorem{lemma}[theorem]{Lemma}
\newtheorem{corollary}[theorem]{Corollary}
\newtheorem{proposition}[theorem]{Proposition}
\theoremstyle{definition}
\newtheorem{definition}[theorem]{Definition}
\newtheorem{remark}[theorem]{Remark}
\newtheorem{example}[theorem]{Example}

% ── Code listings ────────────────────────────────────────────────────────────
\definecolor{codegray}{rgb}{0.96,0.96,0.96}
\definecolor{codeblue}{rgb}{0.1,0.2,0.6}
\definecolor{codegreencomment}{rgb}{0.13,0.55,0.13}

\lstset{
  backgroundcolor=\color{codegray},
  basicstyle=\ttfamily\small,
  keywordstyle=\color{codeblue}\bfseries,
  commentstyle=\color{codegreencomment}\itshape,
  breaklines=true,
  frame=single,
  framerule=0pt,
  rulecolor=\color{gray!30},
  xleftmargin=6pt, xrightmargin=6pt,
}

\lstdefinelanguage{Lean4}{
  keywords={theorem,lemma,def,axiom,by,have,obtain,exact,intro,rw,push_cast,
            ring,decide,linarith,nlinarith,rcases,calc,simp,fin_cases,revert,
            import,where,fun,if,then,else,let,in,match,with,show,apply,refine,
            constructor,left,right,norm_num,norm_cast,exact_mod_cast,
            unfold,push_neg,omega,Set,Finset,open,use,private,rintro,neg_inj,
            linear_combination},
  sensitive=true,
  morecomment=[l]{--},
  morecomment=[s]{/-}{-/},
  morestring=[b]",
  keywordstyle=\color{codeblue}\bfseries,
  commentstyle=\color{codegreencomment}\itshape,
}

% ── Macros ───────────────────────────────────────────────────────────────────
\newcommand{\Z}{\mathbb{Z}}
\newcommand{\Q}{\mathbb{Q}}
\newcommand{\N}{\mathbb{N}}

% ── Header / Footer ──────────────────────────────────────────────────────────
\pagestyle{fancy}
\fancyhf{}
\rhead{\small\thepage}
\lhead{\small Infinitely Many Solutions to $4x^5 + 4y^5 + 11z^5 = 0$}
\renewcommand{\headrulewidth}{0.4pt}

% ── Title ────────────────────────────────────────────────────────────────────
\title{%
  \Large\bfseries Infinitely Many Integer Solutions to\\[6pt]
  \Huge $4x^5 + 4y^5 + 11z^5 = 0$\\[10pt]
  \large Complete Analysis with Lean~4 Formalisation
}
\author{%
  \normalsize Verified in Lean~4 / Mathlib~4 (v4.21.0)\\[2pt]
  \normalsize Repository: \href{https://github.com/JAgbanwa/miscellaneousmathproblems}%
    {\texttt{JAgbanwa/miscellaneousmathproblems}}
}
\date{\normalsize April 2026}

% ════════════════════════════════════════════════════════════════════════════
\begin{document}

\maketitle
\tableofcontents
\newpage

% ════════════════════════════════════════════════════════════════════════════
\section{Statement and Main Result}\label{sec:intro}

\subsection{The equation}

We study the Diophantine equation
\begin{equation}\label{eq:main}
  4x^5 + 4y^5 + 11z^5 = 0, \qquad (x,y,z)\in\Z^3. \tag{$\star$}
\end{equation}

\subsection{Main theorem}

\begin{theorem}[Main result]\label{thm:main}
  The equation $4x^5 + 4y^5 + 11z^5 = 0$ has \textbf{infinitely many} integer solutions.
  In particular, the parametric family
  \begin{equation}\label{eq:family}
    (x,\, y,\, z) = (t,\; -t,\; 0), \qquad t \in \Z,
  \end{equation}
  provides infinitely many distinct integer solutions.
\end{theorem}

\begin{proof}
  Direct verification:
  $4t^5 + 4(-t)^5 + 11\cdot 0^5 = 4t^5 - 4t^5 + 0 = 0$.
  The solutions are distinct because distinct $t$ give different triples.
\end{proof}

% ════════════════════════════════════════════════════════════════════════════
\section{Derivation of the Parametric Family}\label{sec:derivation}

\subsection{Step 1: Setting $z = 0$}

With $z = 0$, equation~\eqref{eq:main} reduces to
\[
  4x^5 + 4y^5 = 0 \implies x^5 = -y^5.
\]
Since the map $t \mapsto t^5$ is strictly monotone on $\Z$ (for $t > s$, factoring
$t^5 - s^5 = (t-s)(t^4 + t^3s + t^2s^2 + ts^3 + s^4) > 0$), we have
\[
  x^5 = -y^5 \iff x^5 = (-y)^5 \iff x = -y.
\]
So the complete set of solutions with $z = 0$ is $\{(t, -t, 0) : t \in \Z\}$.

\subsection{Step 2: Degree-5 homogeneity}

The equation $4x^5 + 4y^5 + 11z^5 = 0$ is \emph{homogeneous of degree~5}: every
monomial has degree exactly 5.  Consequently, if $(a, b, c)$ is any solution, then
for any $t \in \Z$,
\begin{align}
  4(ta)^5 + 4(tb)^5 + 11(tc)^5
  &= t^5\bigl(4a^5 + 4b^5 + 11c^5\bigr) = t^5 \cdot 0 = 0.
\end{align}

The seed solution $(1, -1, 0)$ (which is easily verified directly: $4 - 4 + 0 = 0$)
scales to $(t, -t, 0)$ for all $t \in \Z$, reproducing \eqref{eq:family}.

\begin{remark}
  This is the simplest possible homogeneity structure in this family of problems:
  every variable has the same weight~1, so the equation lives on a standard
  projective variety in $\mathbb{P}^2$.  The rational point
  $(1 : -1 : 0) \in \mathbb{P}^2(\Q)$ corresponds to our parametric family.
\end{remark}

% ════════════════════════════════════════════════════════════════════════════
\section{Modular Analysis}\label{sec:modular}

\subsection{Modulo 2}

Since $4x^5 \equiv 0 \pmod{2}$ and $4y^5 \equiv 0 \pmod{2}$, we have
$11z^5 \equiv z^5 \equiv z \equiv 0 \pmod{2}$, so $z$ must be \textbf{even}.  The
primary family has $z = 0$, which is trivially even. $\checkmark$

\subsection{Modulo 11}

Reducing modulo 11:
\[
  4x^5 + 4y^5 \equiv 0 \pmod{11} \implies x^5 + y^5 \equiv 0 \pmod{11}.
\]
By Fermat's little theorem, $a^{10} \equiv 1 \pmod{11}$ for $\gcd(a,11)=1$, so
$a^5 \equiv \pm 1 \pmod{11}$.  The fifth-power residues modulo~11 are
$\{0, 1, 10\} = \{0, 1, -1\}$.  Thus $x^5 + y^5 \equiv 0 \pmod{11}$ holds when:
\begin{enumerate}
  \item $11 \mid x$ and $11 \mid y$; or
  \item $x^5 \equiv 1$ and $y^5 \equiv -1 \pmod{11}$ (or vice versa).
\end{enumerate}
The primary family satisfies Case~2 when $11 \nmid t$ (since $t^5 + (-t)^5 = 0$) and
Case~1 when $11 \mid t$. $\checkmark$

\subsection{No modular obstruction}

The equation has solutions modulo every prime $p$ (the family $(1,-1,0)$ works for all
$p$ since $4 - 4 = 0$ independently of the residue).  No modular obstruction exists.

% ════════════════════════════════════════════════════════════════════════════
\section{Solutions with $z \neq 0$}\label{sec:z-nonzero}

The classification of all integer solutions, including those with $z \neq 0$, is
equivalent to finding rational points on the affine curve
\[
  C\colon\quad u^5 + v^5 = -\tfrac{11}{4},
  \qquad u = x/z,\; v = y/z \in \Q.
\]
The smooth projective model of $C$ is a Fermat-type quintic of \textbf{genus 6}.
By Faltings' theorem~\cite{faltings83}, $C(\Q)$ is finite.  To determine whether
$C(\Q)$ is empty (i.e., whether no solutions with $z \neq 0$ exist), one would need
a Chabauty--Coleman computation or a $p$-adic descent argument; this is beyond
the scope of the present treatment.

\begin{remark}
  A brute-force search over $|x|, |y| \leq 50$ (with $z$ computed from the
  divisibility condition $11 \mid 4(x^5 + y^5)$) finds \textbf{no solutions with
  $z \neq 0$} within the search range.
\end{remark}

% ════════════════════════════════════════════════════════════════════════════
\section{Explicit Solutions}\label{sec:explicit}

The following table lists solutions from the primary family $(\star)$:

\begin{center}
\begin{tabular}{cccc}
  \toprule
  $x$ & $y$ & $z$ & Check \\
  \midrule
  $0$ & $0$ & $0$ & $0$ \\
  $1$ & $-1$ & $0$ & $4 - 4 + 0 = 0$ \\
  $-1$ & $1$ & $0$ & $-4 + 4 + 0 = 0$ \\
  $2$ & $-2$ & $0$ & $128 - 128 + 0 = 0$ \\
  $-2$ & $2$ & $0$ & $-128 + 128 + 0 = 0$ \\
  $3$ & $-3$ & $0$ & $972 - 972 + 0 = 0$ \\
  $5$ & $-5$ & $0$ & $12500 - 12500 + 0 = 0$ \\
  $10$ & $-10$ & $0$ & $400000 - 400000 + 0 = 0$ \\
  \bottomrule
\end{tabular}
\end{center}

% ════════════════════════════════════════════════════════════════════════════
\section{Lean 4 Formalisation}\label{sec:lean}

The file \texttt{InfinitelyManySolutions.lean} contains a complete, unconditional
Lean~4 / Mathlib~4 formalisation of Theorem~\ref{thm:main}.

\textbf{Sorry count: 0.  Axiom count: 0.}

\subsection{Proof structure}

\begin{lstlisting}[language=Lean4]
import Mathlib

def isolution (x y z : Z) : Prop := 4 * x ^ 5 + 4 * y ^ 5 + 11 * z ^ 5 = 0

-- Degree-5 homogeneity
theorem homogeneity (x y z : Z) (h : isolution x y z) (t : Z) :
    isolution (t * x) (t * y) (t * z) := by
  unfold isolution at *
  linear_combination t ^ 5 * h

-- Primary parametric family (t, -t, 0)
theorem family (t : Z) : isolution t (-t) 0 := by
  unfold isolution; ring

-- Injectivity of n |-> (n, -n, 0) on N
theorem natMap_injective : Function.Injective natMap := by
  intro n1 n2 h
  simp only [natMap, Prod.mk.injEq] at h
  exact_mod_cast h.1

-- Infinitude
theorem solutions_infinite : solutions.Infinite :=
  (Set.infinite_range_of_injective natMap_injective).mono
    (fun _ ⟨n, rfl⟩ => natMap_mem n)
\end{lstlisting}

\subsection{Lean theorem table}

\begin{center}
\begin{tabular}{lll}
  \toprule
  Lean name & Statement & Method \\
  \midrule
  \texttt{homogeneity} & $\forall\, t,\ (tx,ty,tz)$ is a soln if $(x,y,z)$ is & \texttt{linear\_combination} \\
  \texttt{family} & $(t,-t,0)$ is a solution for all $t:\Z$ & \texttt{ring} \\
  \texttt{sol\_1\_neg1\_0} & $(1,-1,0)$ is a solution & \texttt{norm\_num} \\
  \texttt{natMap\_mem} & $(n,-n,0) \in \mathtt{solutions}$ for all $n:\N$ & \texttt{push\_cast} + \texttt{ring} \\
  \texttt{natMap\_injective} & $n \mapsto (n,-n,0)$ is injective & \texttt{exact\_mod\_cast} \\
  \texttt{solutions\_infinite} & solution set is infinite & \texttt{Set.infinite\_range\_of\_injective} \\
  \bottomrule
\end{tabular}
\end{center}

% ════════════════════════════════════════════════════════════════════════════
\section{Proof Status Summary}\label{sec:status}

\begin{center}
\begin{tabular}{ll}
  \toprule
  Component & Status \\
  \midrule
  Seed solution $(1,-1,0)$ & $\checkmark$ Direct substitution \\
  Degree-5 homogeneity & $\checkmark$ Ring identity \\
  Primary family $(t,-t,0)$ verified & $\checkmark$ \texttt{ring} \\
  Infinitely many distinct solutions & $\checkmark$ Injectivity of $n \mapsto n$ on $\N$ \\
  No solutions with $z \neq 0$ (search) & $\checkmark$ Exhaustive check, $|x|,|y| \leq 50$ \\
  No solutions with $z \neq 0$ (proof) & Open --- requires Faltings/Chabauty--Coleman \\
  Lean~4 formalisation (0 sorry) & $\checkmark$ \textbf{Complete} \\
  \bottomrule
\end{tabular}
\end{center}

\begin{thebibliography}{9}
\bibitem{faltings83}
  G.~Faltings,
  \textit{Endlichkeitss\"atze f\"ur abelsche Variet\"aten \"uber Zahlk\"orpern},
  Inventiones Mathematicae \textbf{73} (1983), 349--366.
\end{thebibliography}

\end{document}
